% ===================================================================
%  Generado por Erlen
%  Compílalo en Overleaf, o localmente con dos pasadas de pdflatex.
% ===================================================================
\documentclass[aspectratio=169,11pt]{beamer}
\usetheme{Boadilla}
\usefonttheme{serif}
% Boadilla es el tema de fábrica más limpio y, con serifas, es la puesta en
% escena de una defensa. En pantalla la letra es Termes; en el PDF, la serifa
% de fábrica de LaTeX. Para que coincidan, añade \usepackage{tgtermes}.
\definecolor{marinoPapel}{HTML}{F4F6FA}
\definecolor{marinoTinta}{HTML}{15243D}
\definecolor{marinoAzul}{HTML}{16345E}
\definecolor{marinoLaton}{HTML}{B08A32}
\setbeamercolor{background canvas}{bg=marinoPapel}
\setbeamercolor{normal text}{fg=marinoTinta}
\setbeamercolor{structure}{fg=marinoAzul}
\setbeamercolor{frametitle}{bg=marinoPapel,fg=marinoAzul}
\setbeamercolor{block title}{bg=white,fg=marinoAzul}
\setbeamercolor{block body}{bg=white}
\setbeamercolor{alerted text}{fg=marinoLaton!65!marinoTinta}
\setbeamerfont{frametitle}{series=\bfseries}
% El trozo corto de latón sobre la regla del título, el canto de color de las
% cajas y la barra de sección son de Erlen: Beamer no los trae y no salen
% sin reescribir las plantillas de frametitle y de bloque.
\usepackage[utf8]{inputenc}
\usepackage[T1]{fontenc}
% Español: se carga solo si babel-spanish está instalado (Overleaf lo trae).
\IfFileExists{spanish.ldf}{\usepackage[spanish,es-nodecimaldot]{babel}}{}
\usepackage{amsmath,amssymb}
\usepackage[version=4]{mhchem}
\usepackage{booktabs}
\definecolor{erlenfondo}{HTML}{F4F6FA}
\definecolor{erlenacento}{HTML}{B08A32}
\usepackage{tikz}
\usetikzlibrary{arrows.meta, calc}
\renewcommand{\figurename}{Figura}
\renewcommand{\tablename}{Tabla}
\usepackage{siunitx}
\sisetup{per-mode=symbol, range-units=single}
\setbeamertemplate{navigation symbols}{}
% Algunos temas (Metropolis entre ellos) insertan solos una diapositiva al
% empezar cada sección. Erlen ya emite la suya, así que se desactiva la
% automática: de otro modo cada sección saldría dos veces y la numeración
% del PDF dejaría de coincidir con la de la app.
\AtBeginSection{}
\AtBeginSubsection{}

\title[Defensa de tesis]{Defensa de tesis}
\subtitle{Ejemplo didáctico · datos ilustrativos}
\date{\today}

\begin{document}

% --- diapositiva 1 ---
\begin{frame}[plain]
  \transdissolve[duration=0.25]
  \titlepage
\end{frame}
\note{%
  {\scriptsize\textbf{Tiempo previsto: 0:30 min}}\par\medskip
  Ejemplo didáctico. Sustituye contenido y datos por los de tu trabajo. No atribuir estas cifras a un experimento.\par
}

% --- diapositiva 2 ---
\begin{frame}{La pregunta que guía el trabajo}
  \transdissolve[duration=0.25]
  {\large ¿Puede una modificación del material mejorar su estabilidad?}
  \medskip

  {\small Caso de uso ilustrativo. Define aquí el sistema, la comparación y el alcance.}
\end{frame}
\note{%
  {\scriptsize\textbf{Tiempo previsto: 1:30 min}}\par\medskip
}

% --- diapositiva 3 ---
\begin{frame}{Una estrategia explícita}
  \transdissolve[duration=0.25]
  \begin{figure}
    \centering
    \definecolor{sa0f}{HTML}{CFDCE6}
    \definecolor{sa1f}{HTML}{9DB7CC}
    \definecolor{sa2f}{HTML}{6B93B2}
    \definecolor{sa3f}{HTML}{396E98}
    \definecolor{sat1843819360}{HTML}{15181C}
    \definecolor{sat1226267613}{HTML}{FFFFFF}
    % Con babel-spanish las palabras largas se parten solas y el texto respira mejor.
    \begin{tikzpicture}[x=1cm,y=1cm,font=\footnotesize]
      \draw[fill=sa0f, draw=none] (0.000,0.000) -- (2.226,0.000) -- (2.518,-1.254) -- (2.226,-2.508) -- (0.000,-2.508) -- cycle;
      \draw[fill=sa1f, draw=none] (2.644,0.000) -- (4.870,0.000) -- (5.162,-1.254) -- (4.870,-2.508) -- (2.644,-2.508) -- (2.936,-1.254) -- cycle;
      \draw[fill=sa2f, draw=none] (5.288,0.000) -- (7.514,0.000) -- (7.806,-1.254) -- (7.514,-2.508) -- (5.288,-2.508) -- (5.580,-1.254) -- cycle;
      \draw[fill=sa3f, draw=none] (7.932,0.000) -- (10.157,0.000) -- (10.450,-1.254) -- (10.157,-2.508) -- (7.932,-2.508) -- (8.224,-1.254) -- cycle;
      \node[text=sat1843819360, align=center, text width=1.75cm, inner sep=1pt] at (1.113,-1.254) {\bfseries Definir una comparación controlada};
      \node[text=sat1843819360, align=center, text width=1.75cm, inner sep=1pt] at (3.976,-1.254) {\bfseries Caracterizar ambos materiales};
      \node[text=sat1843819360, align=center, text width=1.75cm, inner sep=1pt] at (6.620,-1.254) {\bfseries Medir estabilidad};
      \node[text=sat1226267613, align=center, text width=1.75cm, inner sep=1pt] at (9.264,-1.254) {\scriptsize \bfseries Evaluar incertidumbre};
    \end{tikzpicture}
  \end{figure}
\end{frame}
\note{%
  {\scriptsize\textbf{Tiempo previsto: 1:30 min}}\par\medskip
  Explica las decisiones y los controles, no solo los nombres de las técnicas.\par
}

% --- diapositiva 4 ---
\begin{frame}{Mostrar la evidencia}
  \transdissolve[duration=0.25]
  \begin{table}
    \centering
    \caption{Ejemplo de organización de resultados; no evidencia experimental.}
    \begin{tabular}{crc}
      \toprule
      \textbf{Condición} & \textbf{Respuesta normalizada} & \textbf{Naturaleza} \\
      \midrule
      Control & 1.00 & Ilustrativa \\
      Modificación A & 1.08 & Ilustrativa \\
      Modificación B & 0.97 & Ilustrativa \\
      \bottomrule
    \end{tabular}
  \end{table}
\end{frame}
\note{%
  {\scriptsize\textbf{Tiempo previsto: 1:30 min}}\par\medskip
  Datos o modelos didácticos: reemplazar antes de presentar resultados propios.\par
}

% --- diapositiva 5 ---
\begin{frame}{Interpretación y límites}
  \transdissolve[duration=0.25]
  \begin{itemize}
    \item Distinguir contribución, observación e hipótesis.
    \item Las diferencias de esta tabla no demuestran significación.
    \item Explicar qué limitaciones cambian la conclusión.
  \end{itemize}
\end{frame}
\note{%
  {\scriptsize\textbf{Tiempo previsto: 1:30 min}}\par\medskip
}

% --- diapositiva 6 ---
\begin{frame}{Antes de compartir}
  \transdissolve[duration=0.25]
  \begin{itemize}
    \item Pregunta e hipótesis explícitas
    \item Métodos suficientes para reproducir la comparación
    \item Datos y referencias propios, comprobados
  \end{itemize}
  \medskip

  {\small Conserva el proyecto JSON y revisa la exportación final.}
\end{frame}
\note{%
  {\scriptsize\textbf{Tiempo previsto: 1:30 min}}\par\medskip
}

\end{document}